% Clase 19 --- El signo de la caída en una bobina (§4.4).
% La regla práctica que hay que retener para los circuitos RL de la clase que
% viene: V_B - V_A = -L dI/dt, con el mismo signo crezca o decrezca la corriente
% (igual que -RI en una resistencia). Los paneles comparan los dos casos con la
% misma referencia de bornes.
\begin{center}
\begin{tikzpicture}[scale=1]

  % =============== (a) corriente creciente ===============
  \begin{scope}
    \draw (0,0) to[short, -*] (0.7,0);
    \draw (0.7,0) to[L] (2.9,0);
    \draw (2.9,0) to[short, *-] (3.6,0);
    \node[figlbl, anchor=north] at (0.7,-0.15) {$A$};
    \node[figlbl, anchor=north] at (2.9,-0.15) {$B$};
    \draw[figvecres, line width=1.2pt] (0.9,0.75) -- (2.1,0.75);
    \node[figlbl, figred, anchor=south] at (1.5,0.8) {$I$};
    \node[figlbl, anchor=north, align=center] at (1.8,-0.85)
      {(a) $\dfrac{dI}{dt}>0$:\\[0.2em] $V_B-V_A<0$};
  \end{scope}

  % =============== (b) corriente decreciente ===============
  \begin{scope}[xshift=6.4cm]
    \draw (0,0) to[short, -*] (0.7,0);
    \draw (0.7,0) to[L] (2.9,0);
    \draw (2.9,0) to[short, *-] (3.6,0);
    \node[figlbl, anchor=north] at (0.7,-0.15) {$A$};
    \node[figlbl, anchor=north] at (2.9,-0.15) {$B$};
    \draw[figvecres, line width=1.2pt] (0.9,0.75) -- (2.1,0.75);
    \node[figlbl, figred, anchor=south] at (1.5,0.8) {$I$};
    \node[figlbl, anchor=north, align=center] at (1.8,-0.85)
      {(b) $\dfrac{dI}{dt}<0$:\\[0.2em] $V_B-V_A>0$};
  \end{scope}

  \node[figlbl, anchor=north, align=center] at (5.0,-2.3)
    {en los dos casos:
     $\boxed{V_B-V_A = -L\,\dfrac{dI}{dt}}$};

\end{tikzpicture}\\[0.5em]
{\small\itshape Ésta es la regla que hace falta tener a mano para los circuitos
$RL$, y conviene enunciarla de una manera que no dependa de si la corriente sube
o baja. Fijado el sentido de $I$ de $A$ hacia $B$, la ``caída'' en la bobina es
$V_B - V_A = -L\,dI/dt$, exactamente igual de mecánica que el $-RI$ de una
resistencia: no hay que decidir nada caso por caso, se pone la fórmula con su
signo y la fórmula se encarga. Si la corriente crece, la expresión da negativa
---la bobina se comporta como algo que consume, oponiéndose---, y si decrece da
positiva, es decir la bobina actúa como una fuente que sostiene la corriente. La
diferencia con la resistencia es de fondo, eso sí: la resistencia disipa
irreversiblemente, mientras que la bobina \emph{devuelve} lo que tomó, porque lo
tenía guardado en el campo magnético. Esa energía almacenada es el tema con el
que sigue la próxima clase.}
\end{center}
